\documentclass[10pt,a4paper]{article}
\usepackage[margin=1.55cm,headheight=13pt]{geometry}
\usepackage[T1]{fontenc}
\usepackage[utf8]{inputenc}
\usepackage{amsmath,amssymb}
\usepackage{graphicx}
\usepackage{booktabs}
\usepackage{array}
\usepackage{tabularx}
\usepackage{multicol}
\usepackage{enumitem}
\usepackage[parfill]{parskip}
\usepackage[dvipsnames]{xcolor}
\usepackage{fancyhdr}
\usepackage{lastpage}
\usepackage[hidelinks]{hyperref}
\usepackage{float}
\setlength{\parskip}{3pt}
\setlength{\parindent}{0pt}
\setlist[itemize]{leftmargin=1.2em, itemsep=1pt, topsep=1pt}
\renewcommand{\arraystretch}{1.05}
\pagestyle{fancy}
\fancyhf{}
\fancyhead[L]{\small AI-ENG-201 -- HW1 Report}
\fancyhead[R]{\small Fall 2026}
\fancyfoot[L]{\small AI Academy, National AI Center}
\fancyfoot[R]{\small Page \thepage\ of \pageref{LastPage}}
\newcommand{\figpath}{figures/}
\begin{document}

\begin{center}
{\small \textsc{AI Academy, National AI Center} \\ AI-ENG-201 -- Machine Learning -- Fall 2026}\\[0.35em]
{\Large\bfseries Homework 1 -- k-Nearest Neighbors Report}\\[-0.1em]
\rule{0.86\textwidth}{0.4pt}
\end{center}

\begin{tabularx}{\textwidth}{@{}lXlX@{}}
\textbf{Name:} & Fatima Alibabayeva & \textbf{Student ID:} & Not provided \\
\textbf{Submitted:} & 5 June & \textbf{Late days used:} & 0 \\
\textbf{Total time:} & $\approx$14 hours & \textbf{Collaborators:} & None \\
\end{tabularx}

\smallskip
\fbox{\begin{minipage}{0.97\textwidth}\small
\textit{I, Fatima Alibabayeva, affirm that the work in this submission is my own. I have not shared or received code, plots, or written analysis for this assignment. Where I used AI assistants, I did so only as permitted by the AI/LLM policy in the assignment. I can explain every line of code and every paragraph of writing in my submission.}\\[3pt]
\textbf{Signed:} Fatima Alibabayeva \hfill \textbf{Date:} 5 June\\[2pt]
\textbf{AI tool usage declaration:} Used AI assistance only for editing, consistency checks, and formatting after drafting; all code behavior and reported metrics were verified by running the submitted scripts.
\end{minipage}}

\section*{Executive summary}
The best single-split validation model is KNN with $k=21$, selected only by validation F1. On validation it obtains accuracy $0.7433$, precision $0.7463$, recall $0.5000$, F1 $0.5988$, and AUC-ROC $0.7169$; on the held-out test set used once at the end it obtains accuracy $0.6730$, F1 $0.5222$, and AUC-ROC $0.7349$. The model clearly beats the majority baseline, whose F1 is $0.0000$ and AUC is $0.5000$. Cross-validation on the train+validation data prefers $k=5$ with mean CV F1 $0.5849$, showing that the single validation split has some selection variance. Standardizing features inside each training fold improves best mean CV F1 from $0.5849$ to $0.7196$ because unscaled fare values dominate Euclidean distance.

\section{Exploratory Data Analysis [20 pts]}
\subsection{Summary statistics [4 pts]}
The Vanderbilt Titanic dataset has $N=1309$ passengers. Numeric summaries are computed before imputation.

\begin{minipage}{0.54\textwidth}
\centering\small
\begin{tabular}{lrrrrr}
\toprule
Feature & Mean & Median & Std & Non-null & Missing\\
\midrule
pclass & 2.295 & 3.000 & 0.838 & 1309 & 0\\
age & 29.881 & 28.000 & 14.413 & 1046 & 263\\
sibsp & 0.499 & 0.000 & 1.042 & 1309 & 0\\
parch & 0.385 & 0.000 & 0.866 & 1309 & 0\\
fare & 33.295 & 14.454 & 51.759 & 1308 & 1\\
\bottomrule
\end{tabular}
\end{minipage}\hfill
\begin{minipage}{0.42\textwidth}
\centering\small
\begin{tabular}{llrr}
\toprule
Feature & Value & Count & \%\\
\midrule
sex & male & 843 & 64.4\\
sex & female & 466 & 35.6\\
embarked & S & 914 & 69.8\\
embarked & C & 270 & 20.6\\
embarked & Q & 123 & 9.4\\
embarked & missing & 2 & 0.2\\
\bottomrule
\end{tabular}
\end{minipage}

Age has 20.1\% missingness and fare has one missing value. Fare is strongly right-skewed: its mean is $33.30$ but its median is only $14.45$.

\subsection{Target distribution and class balance [4 pts]}
\begin{minipage}{0.33\textwidth}
\includegraphics[width=\linewidth]{\figpath target_dist.pdf}
\end{minipage}\hfill
\begin{minipage}{0.63\textwidth}
Class proportions are survived $=0.38$ and did not survive $=0.62$. I consider the dataset \textbf{mildly imbalanced}: the minority class is large enough for standard metrics to be meaningful, but the 62/38 split makes F1 more informative than accuracy because a majority predictor already reaches 61.69\% validation accuracy.
\end{minipage}

\subsection{Three feature-vs-target investigations [4 pts]}
\begin{figure}[H]
\centering
\includegraphics[width=0.31\textwidth]{\figpath sex_survival.pdf}\hfill
\includegraphics[width=0.31\textwidth]{\figpath age_survival.pdf}\hfill
\includegraphics[width=0.31\textwidth]{\figpath fare_survival.pdf}
\caption{Feature-vs-target plots for sex, age, and fare.}
\end{figure}
Sex is highly predictive: female survivors substantially outnumber female non-survivors, while the reverse is true for males. Age alone is weaker because both survival groups have median age about 28. Fare is more informative: survivors have median fare about 26.0, while non-survivors have median fare about 10.5, reflecting the link between class, fare, and survival.

\subsection{Missing-data strategy [4 pts]}
For honest evaluation, imputation values in the experiment script are fit on the training partition only and then applied to validation/test. In cross-validation, imputation is fit inside each training fold. Age and fare use median imputation because both are numeric and fare is skewed; embarked uses mode imputation because it is categorical and has only 2 missing values. Cabin is dropped because 77.5\% of its values are missing, above the 30\% threshold.

\subsection{Encoding and final feature matrix [4 pts]}
Sex is binary encoded as male $\mapsto 1$, female $\mapsto 0$. Embarked is one-hot encoded as \texttt{embarked\_C}, \texttt{embarked\_Q}, and \texttt{embarked\_S} because ports have no natural order. Pclass is kept ordinal because first, second, and third class have a meaningful ranking. The final matrix has $N=1309$ and $p=9$ features: pclass, sex, age, sibsp, parch, fare, and three embarked indicators.

\section{KNN from Scratch [30 pts]}
\subsection{Class API [8 pts]}
\texttt{src/knn.py} implements \texttt{KNN(k=5, metric="euclidean", q=2.0, task="classification", weights="uniform")} with \texttt{fit}, \texttt{predict}, and \texttt{predict\_proba}. Public methods include type hints and docstrings. \texttt{fit} memorizes $X$ and $y$; \texttt{predict} supports classification and regression; \texttt{predict\_proba} returns an $(n_{queries}, n_{classes})$ probability matrix for classification.

\subsection{Vectorized distance computation [8 pts]}
Distances are computed by broadcasting $X_{query}$ from $(n_q,p)$ to $(n_q,1,p)$ and $X_{train}$ from $(N,p)$ to $(1,N,p)$. The subtraction creates a full $(n_q,N,p)$ difference tensor, then Euclidean, Manhattan, or Minkowski distances are reduced along the feature axis. There is no Python loop over training points. Time complexity is $O(n_qNp)$ and memory complexity is $O(n_qNp)$ for the difference tensor plus $O(n_qN)$ for the distance matrix.

\subsection{Sanity check [4 pts]}
The implementation is tested against \texttt{sklearn.neighbors.KNeighborsClassifier} at seed 42. \texttt{pytest} reports \textbf{14 passed, 0 failed}; predictions match exactly and \texttt{predict\_proba} matches within $10^{-9}$.

\subsection{Probability output [5 pts]}
For class $c$, \texttt{predict\_proba} implements
$\hat p(y=c\mid x)=\frac{1}{k}\sum_{i\in N_k(x)}\mathbf{1}\{y^{(i)}=c\}$.
Therefore KNN probabilities lie on the grid $\{0,1/k,2/k,\ldots,1\}$. For AUC-ROC this creates ties: many samples can share exactly the same score. My \texttt{roc\_auc} handles this by sweeping only unique thresholds, computing each tied threshold once, sorting ROC points by FPR, and integrating with the trapezoidal rule.

\subsection{Computational benchmark [5 pts]}
\begin{figure}[H]
\centering
\includegraphics[width=0.57\textwidth]{\figpath benchmark.pdf}
\caption{Prediction time vs. training-set size $N$ on a fixed test set of 100 query points.}
\end{figure}
The empirical log-log scaling exponent is $1.027$, close to the theoretical $O(Np)$ behavior for fixed $p=9$. Small deviations are expected from NumPy overhead and cache effects.

\section{Honest Evaluation [30 pts]}
\subsection{Stratified split [5 pts]}
The split sizes are train $(785,9)$, validation $(261,9)$, and test $(263,9)$. Positive-class proportions are 38.2\% in train, 38.3\% in validation, and 38.0\% in test, so the maximum deviation is only 0.3 percentage points, within the 0.5 pp requirement. The function is verified against \texttt{train\_test\_split(stratify=y)}.

\subsection{Five metrics from scratch [10 pts]}
\texttt{src/metrics.py} implements accuracy, precision, recall, F1, and ROC-AUC without sklearn shortcuts. Precision and recall return 0.0 for zero-denominator edge cases, which makes the majority baseline well-defined. The ROC-AUC implementation explicitly handles KNN score ties using unique thresholds. All metrics are verified against sklearn metric functions at fixed seed with numerical tolerance $10^{-9}$.

\subsection{Hyperparameter sweep [8 pts]}
\begin{figure}[H]
\centering
\includegraphics[width=0.98\textwidth]{\figpath k_sweep.pdf}
\caption{Validation metrics for $k\in\{1,3,5,7,9,11,15,21,31,51\}$.}
\end{figure}
\begin{center}\small
\begin{tabular}{rrrrrr}
\toprule
$k$ & Accuracy & Precision & Recall & F1 & AUC\\
\midrule
1 & 0.6628 & 0.5667 & 0.5100 & 0.5368 & 0.6339\\
3 & 0.6897 & 0.6067 & 0.5400 & 0.5714 & 0.7104\\
5 & 0.6858 & 0.6098 & 0.5000 & 0.5495 & 0.7047\\
7 & 0.7126 & 0.6471 & 0.5500 & 0.5946 & 0.7043\\
21 & 0.7433 & 0.7463 & 0.5000 & \textbf{0.5988} & 0.7169\\
31 & 0.7356 & 0.7246 & 0.5000 & 0.5917 & 0.7286\\
51 & 0.7356 & 0.7460 & 0.4700 & 0.5767 & 0.7261\\
\bottomrule
\end{tabular}
\end{center}
Best $k$ by validation F1 is $k=21$. Recall is highest at smaller $k$, while larger $k$ improves precision and accuracy; F1 chooses the best balance.

\subsection{Baselines [4 pts]}
\begin{center}\small
\begin{tabular}{lrrrrr}
\toprule
Model & Accuracy & Precision & Recall & F1 & AUC-ROC\\
\midrule
Majority baseline & 0.6169 & 0.0000 & 0.0000 & 0.0000 & 0.5000\\
My KNN ($k=21$) & 0.7433 & 0.7463 & 0.5000 & 0.5988 & 0.7169\\
sklearn KNN ($k=21$) & 0.7433 & 0.7463 & 0.5000 & 0.5988 & 0.7172\\
\bottomrule
\end{tabular}
\end{center}
My KNN strongly outperforms the majority baseline and matches sklearn on predictions. The AUC difference of about 0.0003 is due to implementation-level handling of tied probability thresholds, not a meaningful modelling difference.

\subsection{Final test-set evaluation [3 pts]}
The test set was used once after fixing $k=21$ from validation F1.
\begin{center}\small
\begin{tabular}{lrrrrr}
\toprule
Split & Accuracy & Precision & Recall & F1 & AUC-ROC\\
\midrule
Validation & 0.7433 & 0.7463 & 0.5000 & 0.5988 & 0.7169\\
Test & 0.6730 & 0.5875 & 0.4700 & 0.5222 & 0.7349\\
\bottomrule
\end{tabular}
\end{center}
The test F1 is lower than validation by 7.7 pp, mainly due to lower precision on the held-out sample. However, test AUC is slightly higher than validation, so the ranking quality is still consistent; the gap is plausible for a validation/test size around 260 passengers and indicates split variance rather than test-set leakage.

\section{Cross-Validation, Bias-Variance, and Scaling [20 pts]}
\subsection{Stratified K-fold CV [5 pts]}
\texttt{stratified\_kfold} returns 5 folds preserving class proportions. It is verified against \texttt{StratifiedKFold}; every fold has a positive-class rate close to the global 38.2\% rate.

\subsection{CV F1 vs. k [5 pts]}
\begin{figure}[H]
\centering
\includegraphics[width=0.58\textwidth]{\figpath cv_f1_vs_k.pdf}
\caption{5-fold stratified CV mean F1 with $\pm1$ standard deviation band.}
\end{figure}
The best CV choice is $k=5$ with mean F1 $0.5849$ and standard deviation $0.0382$. The curve falls after the mid-range values: $k=31$ has mean F1 $0.5166$ and $k=51$ has mean F1 $0.4891$.

\subsection{CV-best k vs. single-split-best k [3 pts]}
Single validation selects $k=21$, while 5-fold CV selects $k=5$. I trust the CV estimate more for general model selection because it averages over five validation folds instead of relying on one random validation partition. The disagreement is also a useful warning: the single split rewarded high precision at $k=21$, but across folds that choice is less stable than $k=5$.

\subsection{Bias-variance discussion [4 pts]}
\begin{figure}[H]
\centering
\includegraphics[width=0.58\textwidth]{\figpath train_vs_val_f1.pdf}
\caption{Training vs validation F1 across $k$.}
\end{figure}
At $k=1$, training F1 is $0.9511$ while validation F1 is only $0.5368$, a large generalization gap that indicates overfitting. Increasing $k$ smooths the decision boundary and lowers variance: by $k=7$ validation F1 improves to $0.5946$ while training F1 falls to $0.6678$. At very large $k$, the model becomes too smooth; CV F1 drops from $0.5849$ at $k=5$ to $0.4891$ at $k=51$, showing underfitting.

\subsection{Scaling experiment [3 pts]}
\begin{figure}[H]
\centering
\includegraphics[width=0.58\textwidth]{\figpath scaling_effect.pdf}
\caption{Effect of StandardScaler on CV F1; scaler is fit only on the training fold.}
\end{figure}
Without scaling, best mean CV F1 is $0.5849$ at $k=5$. With StandardScaler fit only on each training fold, best mean CV F1 becomes $0.7196$ at $k=5$, a lift of $+13.47$ percentage points. KNN is scale-sensitive because Euclidean distance uses raw numeric ranges: fare ranges up to about 512, while binary variables such as sex and embarked range only from 0 to 1. Without scaling, fare dominates distance; after scaling, sex, pclass, age, and embarked can contribute meaningfully.

\section*{Bonus note}
Weighted KNN is implemented with \texttt{weights="distance"} using weights $1/(d+10^{-8})$, and the equidistant case is tested so it reduces to uniform voting. \texttt{ApproxKNN} is also provided as a BallTree wrapper with the same public API, but the core report grade does not rely on the optional large synthetic benchmark.

\section*{Reproducibility}
All random operations are seeded with 42 unless otherwise stated. Running \texttt{python src/experiments.py} from the project root regenerates the figures and prints the same metric tables; running \texttt{pytest -q} gives 14 passed tests.

\end{document}
